\intro{

\textbf{Actuality.} 




%8. У вступі подається загальна характеристика дисертації, а саме:
%обґрунтування вибору теми дослідження (висвітлюється зв’язок теми дисертації із
%сучасними дослідженнями у відповідній галузі знань шляхом критичного аналізу з визначенням
%сутності наукової проблеми або завдання);
%мета і завдання дослідження відповідно до предмета та об’єкта дослідження;
%методи дослідження (перераховуються використані наукові методи дослідження та
%змістовно відзначається, що саме досліджувалось кожним методом; обґрунтовується вибір
%методів, що забезпечують достовірність отриманих результатів та висновків);
%наукова новизна отриманих результатів (аргументовано, коротко та чітко представляються
%основні наукові положення, які виносяться на захист, із зазначенням відмінності %одержаних
%результатів від відомих раніше);
%особистий внесок здобувача (якщо у дисертації використано ідеї або розробки, що
%належать співавторам, разом з якими здобувачем опубліковано наукові праці, обов’язково
%зазначається конкретний особистий внесок здобувача в такі праці або розробки; здобувач має
%також додати посилання на дисертації співавторів, у яких було використано результати спільних
%робіт);
%апробація матеріалів дисертації (зазначаються назви конференції, конгресу, симпозіуму,
%семінару, школи, місце та дата проведення);
%структура та обсяг дисертації (анонсується структура дисертації, зазначається її загальний
%обсяг).
%За наявності у вступі можуть також вказуватися:
%зв’язок роботи з науковими програмами, планами, темами, грантами - вказується, в рамках
%яких програм, тематичних планів, наукових тематик і грантів, зокрема галузевих, державних та/
%або міжнародних, виконувалося дисертаційне дослідження, із зазначенням номерів державної
%реєстрації науково-дослідних робіт і найменуванням організації, де виконувалася робота;
%практичне значення отриманих результатів - надаються відомості про використання
%результатів досліджень або рекомендації щодо їх практичного використання.
Justification of the choice of the research topic. In the National report on the state and prospects of the development of education in Ukraine, informatization, digitalization, and intellectualization of all spheres of life are indicated as one of the basic transformational processes, and the digitalization of education is a determining factor in its development.
Digital history in science is an interdisciplinary approach to the study of history by digital means. In education and teaching methods, Digital history acts as a new approach to the organization of educational material and the learning process. Using Digital history tools, such as creating interactive maps, processing large amounts of information, recognizing handwritten texts, and creating infographics will make it easier to master the educational material, which in turn expands students' cognitive abilities and facilitates the learning process.
Elements of Digital History are already used in the educational system of Ukraine at all levels. Thus, since 2019, with the beginning of the pandemic, there has been an active use of online courses and distance education, as well as local use of various tools and programs that can be attributed to Digital History.
In this context, Digital History solves several problems of digital transformation at once:

1. Formation and development of digital competencies of participants in the educational process

2. Improvement of the learning process, which relieves students

3. Strengthening ties with humanity, strengthening European integration processes

\textbf{OBJECT OF STUDY}
The educational process of training future professionals in digital history methodologies within higher education institutions.

\textbf{SUBJECT OF STUDY}
The pedagogical methodology for teaching digital history to future professionals, including its theoretical foundations, practical implementation strategies, and effectiveness assessment.

Even though the use of Information and Communication Technologies (ICT) in education is thoroughly presented in the studies of the last decade, we cannot literally use the methods of ICT application, because, in the framework of the processes of globalization and informatization, there is a need for a new, systemic and specialized approach. Digital history is an approach that allows you to combine existing methodological assets with new methods, techniques, and techniques that exist on the border between education and science or are even used in commercial projects. At the same time, for use in general secondary education, this approach needs to be experimentally verified and scientifically understood, that the pain should first be realized with future specialists, as more mobile and independent.
Thus, the relevance of the scientific understanding of the scientific problem the study is due to contradictions between:

• modern needs of society and the insufficient degree of implementation of digitalization and intellectualization in education;

• the importance of the formation of sociocultural competence of future specialists in the conditions of digitalization and the formation of an information society in the absence of a coherent theoretical and experimental justification of this process

• an objective need for scientific methodical support of the process of formation of information competence in future specialists and the insufficient development of the methodology of teaching Digital History in future specialists

Based on this information, the following research question is summoned. 

\textbf{Primary research question}
How can the teaching methodologies of digital history be redesigned to effectively prepare students for professional careers in digital historical fields while maintaining academic rigor and the excellence of historical thinking?

\textbf{Secondary Research Questions}
\begin{enumerate}
  \item How do current theoretical frameworks in digital humanities pedagogy address the specific needs of professional training in digital history?
  \item What baseline digital history competencies do students possess upon entering professional training programs, and how can these be effectively measured?
  \item What institutional and technological conditions are necessary to ensure equitable access to digital history tools in professional training environments?
  \item What ethical frameworks and critical evaluation skills must be embedded in digital history training for future professionals?
\end{enumerate}


\textbf{Primary Research Hypothesis}
"The implementation of a competency-based, industry-integrated digital history teaching methodology will significantly improve both student historical thinking skills and professional readiness outcomes compared to traditional digital history education approaches, while creating measurable value for partner professional institutions through enhanced graduate preparedness and collaborative innovation."

Therefore, the need to resolve the identified contradictions, the insufficient level of theoretical and practical development of the specified problem, the lack of its systematic scientific analysis, and its relevance and novelty determined the choice of the topic of the dissertation work "Methodology of teaching Digital History of future specialists".
Connection of work with scientific programs, plans, topics. The dissertation work was carried out by the plans of the scientific research work of the Department of Pedagogy of the Kryvyi Rih State Pedagogical University and the topic of the complex scientific study "Methodology of teaching Digital History to future specialists" (State registration number 0120U101116 dated 26.02.2020).
The purpose of the study is to theoretically substantiate, develop, and experimentally test the methods of teaching digital history for future specialists.
The object of research is the process of teaching digital history to future specialists
The subject of the study is the method of teaching digital history in institutions of higher education.

.
}